\section{The Boolean-Kernel Basis: Complete Algebraic Theory}\label{sec:kbalgebra}

We develop the Boolean-kernel representation with all indices, maps, fields, and degree conventions explicit. Each result is split into elementary steps suitable for future Lean~4 certification. The original KBFold paper proves the kernel basis, Kronecker transform, low-degree criterion, and fold dictionary~\cite{MkhidaKBFold2026}; we reconstruct those results with complete proofs, then derive algebraic consequences directly relevant to degree-aware masking and bounded-query simulation. Throughout this chapter, $\F$ is an arbitrary field unless a lemma explicitly assumes a finite field or specifies the characteristic.

\subsection{Bit indexing, kernels, and coefficients}\label{sec:kbalgebra:bits}

The statements in this section do not require a finite field or an assumption on the characteristic. They rederive the structural facts of the original KBFold basis~\cite{MkhidaKBFold2026} with complete proofs and explicit indexing conventions.

\begin{definition}[Bit sets and binary exponents]\label{def:kbalgebra:bits}
Let $\F$ be an arbitrary field and let $m\geq0$ be an integer. Write $J_m:=\{0,\ldots,m-1\}$, $I_m:=\{0,1\}^m$, and $N:=2^m$. For $S\subseteq J_m$, define
\[
 e(S):=\sum_{i\in S}2^i\in\{0,\ldots,N-1\}.
\]
For $b\in I_m$, set $\operatorname{supp}(b):=\{i:b_i=1\}$, $Z(b):=J_m\setminus\operatorname{supp}(b)$, and $\operatorname{wt}(b):=|\operatorname{supp}(b)|$. The empty set $J_0$ and empty tuple $I_0=\{()\}$ have their usual meanings.
\end{definition}

\begin{lemma}[Binary-subset bijection]\label{lem:kbalgebra:binary-bijection}
The function $e:\mathcal P(J_m)\to\{0,\ldots,N-1\}$ is a bijection. Its inverse sends the binary expansion $a=\sum_{i=0}^{m-1}a_i2^i$ to $\{i:a_i=1\}$.
\end{lemma}
\begin{proof}
For every $S$, $0\leq e(S)\leq\sum_{i=0}^{m-1}2^i=2^m-1$. If $e(S)=e(T)$ and $S\neq T$, choose the largest $j$ in the symmetric difference. The contribution of $2^j$ to $e(S)-e(T)$ has absolute value $2^j$, whereas the sum of all contributions from smaller indices has absolute value at most $\sum_{i=0}^{j-1}2^i=2^j-1$. Equality is impossible. Thus $e$ is injective. Domain and codomain both have cardinality $2^m$, so it is bijective. The inverse formula follows directly from the definition of $e$.
\end{proof}

\begin{definition}[General kernel and Boolean kernel]\label{def:kbalgebra:kernel}
For $y=(y_i)_{i\in J_m}\in\F^m$ define
\[
 K_y^{(m)}(X):=\prod_{i=0}^{m-1}(X^{2^i}+y_i)\in\F[X].
\]
For $b\in I_m\subseteq\F^m$ write $K_b^{(m)}:=K_y^{(m)}|_{y=b}$. We put $K_{()}^{(0)}=1$. For a table $f:I_m\to\F$, set
\[
 U_f^{(m)}(X):=\sum_{b\in I_m}f(b)K_b^{(m)}(X).
\]
\end{definition}

\begin{lemma}[Explicit product expansion]\label{lem:kbalgebra:product-expansion}
For every $y\in\F^m$,
\begin{equation}\label{eq:kbalgebra:expansion}
 K_y^{(m)}(X)=\sum_{S\subseteq J_m}\left(\prod_{i\in J_m\setminus S}y_i\right)X^{e(S)}.
\end{equation}
\end{lemma}
\begin{proof}
Apply distributivity to the product defining $K_y^{(m)}$. For each factor indexed by $i$, choose either the monomial $X^{2^i}$ or the constant $y_i$. Let $S$ be the set of indices where the monomial is chosen. The resulting term is exactly $\left(\prod_{i\notin S}y_i\right)X^{\sum_{i\in S}2^i}$. Every choice corresponds to precisely one subset $S$, so summing the terms gives~\eqref{eq:kbalgebra:expansion}.
\end{proof}

\begin{lemma}[Coefficient formula]\label{lem:kbalgebra:coefficient-rule}
For $S\subseteq J_m$ and $y\in\F^m$,
\[
 [X^{e(S)}]K_y^{(m)}=\prod_{i\notin S}y_i.
\]
For Boolean $b\in I_m$, this coefficient equals $1$ precisely when $Z(b)\subseteq S$ and equals $0$ otherwise.
\end{lemma}
\begin{proof}
Lemma~\ref{lem:kbalgebra:binary-bijection} shows that the exponents $e(T)$ in~\eqref{eq:kbalgebra:expansion} are pairwise distinct. Thus the only term of exponent $e(S)$ is the one indexed by $S$, proving the first formula. If $y=b$ is Boolean, the product equals $1$ if all factors indexed by $J_m\setminus S$ are $1$, and is zero if any of those factors is $0$. The condition is $(J_m\setminus S)\cap Z(b)=\varnothing$, which is equivalent to $Z(b)\subseteq S$.
\end{proof}

\begin{corollary}[Degree, leading term, and constant term]\label{cor:kbalgebra:monic}
For each $b\in I_m$, $K_b^{(m)}$ is monic of degree $N-1$ and
\[
 [X^0]K_b^{(m)}=\prod_{i\in J_m}b_i,
 \qquad [X^{N-1}]K_b^{(m)}=1.
\]
The same degree and leading coefficient assertion holds for $K_y^{(m)}$ with arbitrary $y\in\F^m$.
\end{corollary}
\begin{proof}
The greatest exponent occurring in~\eqref{eq:kbalgebra:expansion} is $e(J_m)=N-1$, and its coefficient is the empty product $1$. The exponent $e(\varnothing)=0$ has coefficient $\prod_iy_i$. These statements hold also for $m=0$, where $K_{()}^{(0)}=1$.
\end{proof}

\begin{lemma}[Block factorization]\label{lem:kbalgebra:block}
For integers $0\leq t\leq m$, $c\in I_t$, and $h\in I_{m-t}$, write $(c,h)$ for the concatenation of the $t$ low-index bits of $c$ with the $m-t$ high-index bits of $h$. Then
\begin{equation}\label{eq:kbalgebra:block}
 K_{(c,h)}^{(m)}(X)=K_c^{(t)}(X)\,K_h^{(m-t)}(X^{2^t}).
\end{equation}
\end{lemma}
\begin{proof}
Split the defining product into factors with indices $i<t$ and $i\geq t$. The low-index product is $K_c^{(t)}(X)$. In the remaining product put $i=t+j$, so that $X^{2^i}=(X^{2^t})^{2^j}$. The high-index product is therefore $K_h^{(m-t)}(X^{2^t})$. Multiply the two expressions.
\end{proof}

\begin{corollary}[Lowest-bit kernel recursion]\label{cor:kbalgebra:recursion}
For $m\geq1$ and $b'\in I_{m-1}$,
\[
 K_{(0,b')}^{(m)}(X)=XK_{b'}^{(m-1)}(X^2),\qquad
 K_{(1,b')}^{(m)}(X)=(X+1)K_{b'}^{(m-1)}(X^2).
\]
\end{corollary}
\begin{proof}
Take $t=1$ in Lemma~\ref{lem:kbalgebra:block} and note that $K_0^{(1)}(X)=X$ and $K_1^{(1)}(X)=X+1$.
\end{proof}

\subsection{Kronecker transform and the basis theorem}\label{sec:kbalgebra:matrix}

\begin{definition}[Coefficient vectors and kernel coordinates]\label{def:kbalgebra:coords}
For $P\in\F[X]_{<N}$ define its monomial coefficient vector $a\in\F^{I_m}$ by $a(b):=[X^{e(\operatorname{supp}(b))}]P$. For any table $f\in\F^{I_m}$, its encoded polynomial is $U_f^{(m)}$. Define the $N\times N$ matrix $\mathbf K_m$, with rows indexed by $a\in I_m$ and columns by $b\in I_m$, by
\[
 \mathbf K_m[a,b]:=[X^{e(\operatorname{supp}(a))}]K_b^{(m)}.
\]
\end{definition}

\begin{definition}[Bitwise tensor product of matrices]\label{def:kbalgebra:tensor}
For matrices $A_i\in\operatorname{Mat}_{2\times2}(\F)$ indexed by $i\in J_m$, define $\bigotimes_{i=0}^{m-1}A_i$ to be the $I_m\times I_m$ matrix with entries
\[
 \left(\bigotimes_{i=0}^{m-1}A_i\right)[a,b]
 :=\prod_{i=0}^{m-1}A_i[a_i,b_i].
\]
For $m=0$ this is the $1\times1$ identity matrix.
\end{definition}

\begin{lemma}[Tensor multiplication rule]\label{lem:kbalgebra:mixed-product}
For each $i\in J_m$ let $A_i,B_i$ be $2\times2$ matrices over $\F$. Then
\[
 \left(\bigotimes_i A_i\right)\left(\bigotimes_i B_i\right)
 =\bigotimes_i(A_iB_i).
\]
\end{lemma}
\begin{proof}
Fix $a,b\in I_m$. The $[a,b]$ entry of the matrix product on the left is
\[
 \sum_{c\in I_m}\prod_{i\in J_m}A_i[a_i,c_i]B_i[c_i,b_i].
\]
The finite sum ranges over all independent binary choices $c_i\in\{0,1\}$. Distributivity therefore rewrites it as
\[
 \prod_{i\in J_m}\left(\sum_{u\in\{0,1\}}A_i[a_i,u]B_i[u,b_i]\right)
 =\prod_{i\in J_m}(A_iB_i)[a_i,b_i],
\]
which is the $[a,b]$ entry of the right-hand side. Equality of all entries proves the identity. For $m=0$, both sides are the $1\times1$ identity.
\end{proof}

\begin{lemma}[One-bit coefficient transform]\label{lem:kbalgebra:one-bit}
Put
\[
 C=\begin{pmatrix}0&1\\1&1\end{pmatrix},\qquad
 D=\begin{pmatrix}-1&1\\1&0\end{pmatrix}.
\]
Then $CD=DC=I_2$. The vector of monomial coefficients of $f_0K_0^{(1)}+f_1K_1^{(1)}$ is $C(f_0,f_1)^{\mathsf T}=(f_1,f_0+f_1)^{\mathsf T}$, and its inverse is $(f_0,f_1)=(a_1-a_0,a_0)$.
\end{lemma}
\begin{proof}
Since $K_0^{(1)}=X$ and $K_1^{(1)}=X+1$, we have $f_0K_0^{(1)}+f_1K_1^{(1)}=f_1+(f_0+f_1)X$, giving the formula for $C$. Direct matrix multiplication shows
$CD=\left(\begin{smallmatrix}1&0\\0&1\end{smallmatrix}\right)=DC$ over any field, since $-1+1=0$. Therefore $D(a_0,a_1)^{\mathsf T}=(a_1-a_0,a_0)^{\mathsf T}$ is the inverse coefficient map.
\end{proof}

\begin{proposition}[Kronecker coefficient matrix]\label{prop:kbalgebra:kronecker}
With $C$ as in Lemma~\ref{lem:kbalgebra:one-bit},
\[
 \boxed{\mathbf K_m=C^{\otimes m}.}
\]
\end{proposition}
\begin{proof}
The coefficient formula in Lemma~\ref{lem:kbalgebra:coefficient-rule} gives
\[
 \mathbf K_m[a,b]=\prod_{i:a_i=0}b_i.
\]
But $C[0,b_i]=b_i$ and $C[1,b_i]=1$, so the same quantity equals $\prod_{i\in J_m}C[a_i,b_i]$, which is the defining entry of $C^{\otimes m}$. This proves equality entry by entry, including $m=0$.
\end{proof}

\begin{theorem}[Boolean-kernel basis and inverse transform]\label{thm:kbalgebra:basis}
For every field $\F$ and $m\geq0$, the family $\{K_b^{(m)}:b\in I_m\}$ is an $\F$-basis of $\F[X]_{<N}$. Its coefficient matrix and inverse are
\[
 \mathbf K_m=C^{\otimes m},\qquad
 \mathbf K_m^{-1}=D^{\otimes m}.
\]
\end{theorem}
\begin{proof}
By Lemma~\ref{lem:kbalgebra:one-bit}, $CD=DC=I_2$. Applying the mixed-product rule in Lemma~\ref{lem:kbalgebra:mixed-product} to $m$ identical factors gives
\[
 C^{\otimes m}D^{\otimes m}=(CD)^{\otimes m}=I_{I_m},
 \quad D^{\otimes m}C^{\otimes m}=I_{I_m}.
\]
Thus the coefficient matrix is invertible. Its $N$ columns are the monomial coefficient vectors of the $N$ kernel polynomials. Invertibility implies linear independence of these column vectors. The monomials $1,X,\ldots,X^{N-1}$ form an $N$-element basis of $\F[X]_{<N}$, so its dimension is $N$. Therefore an independent set of $N$ kernel polynomials is a basis. No restriction on the characteristic was used.
\end{proof}

\begin{corollary}[Unique kernel coordinates]\label{cor:kbalgebra:unique-coordinates}
For each $P\in\F[X]_{<N}$ there exists a unique $f:I_m\to\F$ with $P=U_f^{(m)}$, and if $a$ is its monomial coefficient vector, then
\[
 a=\mathbf K_m f,\qquad f=\mathbf K_m^{-1}a.
\]
\end{corollary}
\begin{proof}
The existence and uniqueness are the defining properties of a basis in Theorem~\ref{thm:kbalgebra:basis}. Passing to monomial coefficient vectors gives $a=\mathbf K_m f$. Applying its explicit inverse gives the second equality.
\end{proof}

\subsection{M\"obius inversion, constants, and fast transforms}\label{sec:kbalgebra:transforms}

\begin{lemma}[Forward transform in subset form]\label{lem:kbalgebra:forward}
Let $f:I_m\to\F$, set $P:=U_f^{(m)}$, and let $a_S:=[X^{e(S)}]P$ for $S\subseteq J_m$. Then
\begin{equation}\label{eq:kbalgebra:forward}
 \boxed{a_S=\sum_{\substack{b\in I_m\\Z(b)\subseteq S}}f(b).}
\end{equation}
\end{lemma}
\begin{proof}
By Definition~\ref{def:kbalgebra:kernel}, $P=\sum_bf(b)K_b^{(m)}$. Coefficient extraction is $\F$-linear. Thus
\[
 a_S=\sum_b f(b)[X^{e(S)}]K_b^{(m)}.
\]
Lemma~\ref{lem:kbalgebra:coefficient-rule} says the bracket is $1$ precisely when $Z(b)\subseteq S$ and $0$ otherwise. Removing the zero summands gives~\eqref{eq:kbalgebra:forward}.
\end{proof}

\begin{lemma}[Monomial expansion in the kernel basis]\label{lem:kbalgebra:monomial-inverse}
For every $S\subseteq J_m$,
\begin{equation}\label{eq:kbalgebra:monomial-inverse}
 \boxed{X^{e(S)}=\sum_{\substack{b\in I_m\\S\subseteq Z(b)}}
 (-1)^{|Z(b)\setminus S|}K_b^{(m)}(X).}
\end{equation}
\end{lemma}
\begin{proof}
For each $i\in S$ use $X^{2^i}=X^{2^i}+0$. For each $i\notin S$, use the identity
\[
 1=(X^{2^i}+1)-(X^{2^i}+0).
\]
Multiplying these identities gives
\[
 X^{e(S)}=\prod_{i\in S}(X^{2^i}+0)
 \prod_{i\notin S}\left[(X^{2^i}+1)-(X^{2^i}+0)\right].
\]
Expand the second product by distributivity. Each term is indexed by a Boolean vector $b$ with $b_i=0$ for all $i\in S$; equivalently $S\subseteq Z(b)$. Its factors, before their signs, multiply to $K_b^{(m)}$. A negative sign is incurred for every index $i\notin S$ where the term $-(X^{2^i}+0)$ was chosen, i.e. precisely for the indices in $Z(b)\setminus S$. Its coefficient is therefore $(-1)^{|Z(b)\setminus S|}$. Summing these terms proves the equality. This calculation works in every characteristic; in characteristic two all signs reduce to $+1$.
\end{proof}

\begin{proposition}[Inverse M\"obius formula]\label{prop:kbalgebra:inverse}
Let $P(X)=\sum_{S\subseteq J_m}a_S X^{e(S)}$ be arbitrary of degree $<N$. Its unique kernel coordinates $f$ satisfy
\begin{equation}\label{eq:kbalgebra:inverse}
 \boxed{f(b)=\sum_{S\subseteq Z(b)}(-1)^{|Z(b)\setminus S|}a_S.}
\end{equation}
\end{proposition}
\begin{proof}
Substitute the monomial expansion~\eqref{eq:kbalgebra:monomial-inverse} into $P=\sum_Sa_SX^{e(S)}$ and interchange the two finite sums:
\[
 P=\sum_{b\in I_m}\left(\sum_{S\subseteq Z(b)}(-1)^{|Z(b)\setminus S|}a_S\right)K_b^{(m)}.
\]
By uniqueness of kernel coordinates (Corollary~\ref{cor:kbalgebra:unique-coordinates}), the coefficient in parentheses equals $f(b)$.
\end{proof}

\begin{lemma}[One-coordinate transform algorithms]\label{lem:kbalgebra:pair-algorithm}
For a fixed coordinate $i$, apply to every pair of table entries with equal other bits the update
\[
 (x_0,x_1)\longmapsto(x_1,x_0+x_1).
\]
Composing these updates over all $i\in J_m$ sends a kernel-coordinate table $f$ to the monomial coefficient table $a$ of $U_f^{(m)}$. Reversing each update as
\[
 (y_0,y_1)\longmapsto(y_1-y_0,y_0)
\]
recovers $f$ from $a$.
\end{lemma}
\begin{proof}
By Lemma~\ref{lem:kbalgebra:one-bit}, the forward pair update is multiplication by $C$ and the reverse is multiplication by $D=C^{-1}$. When performed on coordinate $i$ while fixing the other bits, the update is multiplication by the tensor matrix that has $C$ at coordinate $i$ and $I_2$ elsewhere. The mixed-product rule of Lemma~\ref{lem:kbalgebra:mixed-product} shows that multiplying these matrices once for each coordinate yields $C^{\otimes m}=\mathbf K_m$. The reverse operations similarly yield $D^{\otimes m}=\mathbf K_m^{-1}$, proving both claims.
\end{proof}

\begin{corollary}[Exact transform cost]\label{cor:kbalgebra:transform-cost}
For $m\geq1$, either transform in Lemma~\ref{lem:kbalgebra:pair-algorithm} takes exactly $m2^{m-1}=mN/2$ additions or subtractions in $\F$ and no nontrivial field multiplications. For $m=0$, both transforms are the identity.
\end{corollary}
\begin{proof}
For each coordinate $i$, there are exactly $2^{m-1}=N/2$ disjoint pairs, since the other $m-1$ bits can be chosen freely. Each forward pair operation copies one entry and computes one sum; each reverse pair copies one entry and computes one difference. Thus each coordinate costs $N/2$ additions or subtractions, and there are $m$ coordinates. Multiplication by $0$, $1$, or $-1$ is not counted as a nontrivial field multiplication. When $m=0$ there are no coordinate updates.
\end{proof}

\begin{corollary}[Highest and constant coefficients]\label{cor:kbalgebra:extreme-coeffs}
For every table $f:I_m\to\F$,
\[
 [X^{N-1}]U_f^{(m)}=\sum_{b\in I_m}f(b),\qquad
 [X^0]U_f^{(m)}=f(1,\ldots,1).
\]
\end{corollary}
\begin{proof}
In Lemma~\ref{lem:kbalgebra:forward}, take $S=J_m$; then every $Z(b)$ is a subset of $S$, so all $f(b)$ appear. For the constant term take $S=\varnothing$; the condition $Z(b)\subseteq\varnothing$ forces $Z(b)=\varnothing$, hence $b=(1,\ldots,1)$, leaving only this term.
\end{proof}

\begin{proposition}[Kernel coordinates of a constant]\label{prop:kbalgebra:constant}
For every $c\in\F$,
\begin{equation}\label{eq:kbalgebra:constant}
 \boxed{c=\sum_{b\in I_m}c(-1)^{m-\operatorname{wt}(b)}K_b^{(m)}(X).}
\end{equation}
\end{proposition}
\begin{proof}
Apply Lemma~\ref{lem:kbalgebra:monomial-inverse} with $S=\varnothing$ to $X^{e(\varnothing)}=X^0=1$. Its sum runs over all $b$, and its sign is $(-1)^{|Z(b)|}=(-1)^{m-\operatorname{wt}(b)}$. Multiply the resulting polynomial identity by $c$.
\end{proof}

\begin{proposition}[Kernel partition product]\label{prop:kbalgebra:partition}
For every $m\geq0$ and every field $\F$,
\begin{equation}\label{eq:kbalgebra:partition}
 \boxed{\sum_{b\in I_m}K_b^{(m)}(X)=\prod_{i=0}^{m-1}(2X^{2^i}+1).}
\end{equation}
If $\operatorname{char}\F=2$, the left side equals $1$; if $\operatorname{char}\F\neq2$ and $m\geq1$, the right side is a nonconstant polynomial of degree $N-1$.
\end{proposition}
\begin{proof}
Since $b_i$ ranges independently over $0$ and $1$, repeated distributivity factors the sum as
\[
 \sum_{(b_i)_i}\prod_i(X^{2^i}+b_i)
 =\prod_i\bigl((X^{2^i}+0)+(X^{2^i}+1)\bigr)
 =\prod_i(2X^{2^i}+1).
\]
In characteristic two, every factor is $1$. Otherwise every factor has leading term $2X^{2^i}$ with nonzero coefficient, so their product has degree $\sum_i2^i=N-1>0$. For $m=0$ both sides are $1$.
\end{proof}

\subsection{Degree constraints in kernel coordinates}\label{sec:kbalgebra:lowdeg}

\begin{definition}[Coefficient blocks]\label{def:kbalgebra:blocks}
Fix integers $0\leq t\leq m$ and set $B:=2^t$, $H:=2^{m-t}$, so $N=BH$. For any $P\in\F[X]_{<N}$, its $B$-coefficient blocks are the polynomials $P_j\in\F[X]_{<B}$, $0\leq j<H$, uniquely determined by
\[
 P(X)=\sum_{j=0}^{H-1}X^{Bj}P_j(X).
\]
\end{definition}

\begin{lemma}[Existence and uniqueness of coefficient blocks]\label{lem:kbalgebra:block-unique}
The decomposition in Definition~\ref{def:kbalgebra:blocks} exists and is unique for every $P\in\F[X]_{<N}$.
\end{lemma}
\begin{proof}
Write $P(X)=\sum_{a=0}^{N-1}p_aX^a$. Every $a$ has a unique quotient--remainder decomposition $a=Bj+i$ with $0\leq i<B$ and $0\leq j<H$. Define $P_j(X):=\sum_{i=0}^{B-1}p_{Bj+i}X^i$. Then
$\sum_jX^{Bj}P_j(X)=\sum_{j,i}p_{Bj+i}X^{Bj+i}=P(X)$. Conversely, comparing coefficients of the pairwise distinct powers $X^{Bj+i}$ determines all coefficients of every $P_j$, proving uniqueness.
\end{proof}

\begin{lemma}[Block degree criterion]\label{lem:kbalgebra:degree-blocks}
With the preceding notation, $\deg P<B$ (including $P=0$) if and only if $P_j=0$ for all $1\leq j<H$.
\end{lemma}
\begin{proof}
The terms of $X^{Bj}P_j$ have degrees in the interval $[Bj,B(j+1)-1]$. These intervals are pairwise disjoint. If $P_j\neq0$ for some $j\geq1$, the polynomial $P$ has a nonzero coefficient of degree at least $B$, so $\deg P\not<B$. Conversely, if all blocks with $j\geq1$ vanish, then $P=P_0$ with $\deg P_0<B$ by definition.
\end{proof}

\begin{definition}[High-bit kernel polynomials]\label{def:kbalgebra:high-block}
For a table $f:I_m\to\F$ and each $c\in I_t$, define
\[
 Y_c(Z):=\sum_{h\in I_{m-t}} f(c,h)K_h^{(m-t)}(Z)\in\F[Z]_{<H}.
\]
If $Y_c(Z)=\sum_{j=0}^{H-1}y_{c,j}Z^j$, define
\[
 A_j(X):=\sum_{c\in I_t}y_{c,j}K_c^{(t)}(X)\in\F[X]_{<B}.
\]
\end{definition}

\begin{lemma}[Kernel blocks equal monomial blocks]\label{lem:kbalgebra:kernel-block-decomposition}
For the polynomials of Definition~\ref{def:kbalgebra:high-block},
\[
 U_f^{(m)}(X)=\sum_{c\in I_t}K_c^{(t)}(X)Y_c(X^{B})
 =\sum_{j=0}^{H-1}X^{Bj}A_j(X).
\]
\end{lemma}
\begin{proof}
Split the finite sum defining $U_f$ according to the unique low/high decomposition $b=(c,h)$. Lemma~\ref{lem:kbalgebra:block} gives
\[
 U_f=\sum_c\sum_hf(c,h)K_c^{(t)}(X)K_h^{(m-t)}(X^B)
 =\sum_cK_c^{(t)}(X)Y_c(X^B).
\]
Substitute $Y_c(X^B)=\sum_jy_{c,j}X^{Bj}$, interchange the two finite sums, and use the definition of $A_j$. By Corollary~\ref{cor:kbalgebra:monic}, each $K_c^{(t)}$ has degree $B-1$, so $A_j\in\F[X]_{<B}$ and the displayed sum is precisely the unique block decomposition.
\end{proof}

\begin{lemma}[Low degree iff each high kernel polynomial is constant]\label{lem:kbalgebra:high-constant}
With the preceding notation,
\[
 \deg U_f^{(m)}<B
 \quad\Longleftrightarrow\quad
 \text{every }Y_c(Z)\text{ is constant.}
\]
\end{lemma}
\begin{proof}
By Lemma~\ref{lem:kbalgebra:degree-blocks} and Lemma~\ref{lem:kbalgebra:kernel-block-decomposition}, $\deg U_f<B$ is equivalent to $A_j=0$ for every $j\geq1$. By the basis theorem (Theorem~\ref{thm:kbalgebra:basis}) with $t$ variables, each $A_j=\sum_cy_{c,j}K_c^{(t)}$ vanishes iff every $y_{c,j}=0$. Thus all $A_j$ with $j\geq1$ vanish iff all coefficients of all $Y_c$ of positive degree vanish, exactly when every $Y_c$ is constant.
\end{proof}

\begin{theorem}[Complete low-degree characterization]\label{thm:kbalgebra:low-degree}
Let $0\leq t\leq m$ and let $f:I_m\to\F$. Put\par\noindent\[U_f^{(m)}=\sum_{b\in I_m}f(b)K_b^{(m)}.\] Then
\begin{equation}\label{eq:kbalgebra:lowdeg}
 \boxed{\deg U_f^{(m)}<2^t}
\end{equation}
if and only if there exists a unique table $\varphi:I_t\to\F$ satisfying
\begin{equation}\label{eq:kbalgebra:lowdeg-pattern}
 \boxed{f(c,h)=(-1)^{(m-t)-\operatorname{wt}(h)}\varphi(c)
 \qquad(c\in I_t,\ h\in I_{m-t}).}
\end{equation}
In this case,
\[
 U_f^{(m)}=U_\varphi^{(t)},
 \qquad \varphi(c)=f(c,1,\ldots,1).
\]
\end{theorem}
\begin{proof}
Put $B=2^t$, $H=2^{m-t}$ and use Definition~\ref{def:kbalgebra:high-block}. By Lemma~\ref{lem:kbalgebra:high-constant}, $\deg U_f<B$ holds if and only if for every $c$, $Y_c(Z)=\varphi(c)$ for some scalar $\varphi(c)$. Proposition~\ref{prop:kbalgebra:constant}, applied with $m-t$ variables, gives the kernel expansion
\[
 \varphi(c)=\sum_{h\in I_{m-t}}
 (-1)^{(m-t)-\operatorname{wt}(h)}\varphi(c)K_h^{(m-t)}(Z).
\]
The definition $Y_c=\sum_hf(c,h)K_h^{(m-t)}$ and uniqueness of kernel coordinates (Corollary~\ref{cor:kbalgebra:unique-coordinates}) imply that $Y_c=\varphi(c)$ iff the coefficients satisfy~\eqref{eq:kbalgebra:lowdeg-pattern} for every $h$. This proves the equivalence. The bit vector $h=(1,\ldots,1)$ has weight $m-t$, so its sign in~\eqref{eq:kbalgebra:lowdeg-pattern} is $+1$. Hence $\varphi(c)=f(c,1,\ldots,1)$ is forced and unique. Finally, substituting $Y_c=\varphi(c)$ into Lemma~\ref{lem:kbalgebra:kernel-block-decomposition} gives
$U_f^{(m)}=\sum_cK_c^{(t)}\varphi(c)=U_\varphi^{(t)}$.
\end{proof}

\begin{corollary}[Dimension of a kernel low-degree subspace]\label{cor:kbalgebra:dimension}
The set of coefficient tables
\[
 \mathcal H_t:=\{f\in\F^{I_m}:\deg U_f^{(m)}<2^t\}
\]
is an $\F$-linear subspace of dimension $2^t$.
\end{corollary}
\begin{proof}
By Theorem~\ref{thm:kbalgebra:low-degree}, $\mathcal H_t$ consists exactly of the tables obtained by applying the linear map
$\varphi\mapsto((c,h)\mapsto(-1)^{m-t-\operatorname{wt}(h)}\varphi(c))$
to $\F^{I_t}$. This map is injective because $\varphi(c)$ is recovered from $f(c,1,\ldots,1)$. Its image is $\mathcal H_t$. Since $\F^{I_t}$ has dimension $|I_t|=2^t$, the same holds for $\mathcal H_t$.
\end{proof}

\begin{remark}[Characteristic-two simplification]\label{rem:kbalgebra:lowdeg-char2}
When $\operatorname{char}\F=2$, every sign in~\eqref{eq:kbalgebra:lowdeg-pattern} equals $+1$, so low-degree kernel tables are constant along each high-bit fiber: $f(c,h)=\varphi(c)$. This simplification fails in general over the odd-characteristic fields used by the original KBFold protocol.
\end{remark}

\subsection{Evaluation on a curve in multilinear coordinates}\label{sec:kbalgebra:evaluation}

\begin{definition}[Multilinear equality weights]\label{def:kbalgebra:eq}
For $b\in I_m$ and $z=(z_i)_{i\in J_m}\in\F^m$, define
\[
 \lambda_b(z):=\prod_{i=0}^{m-1}\bigl((1-b_i)(1-z_i)+b_i z_i\bigr),
 \qquad \widetilde f(z):=\sum_{b\in I_m}f(b)\lambda_b(z).
\]
For $m=0$, $\lambda_{()}(())=1$ and $\widetilde f(())=f(())$.
\end{definition}

\begin{lemma}[Kernel evaluation in factored form]\label{lem:kbalgebra:point-eval}
For $b\in I_m$ and $\zeta\in\F$, write $t_i:=\zeta^{2^i}$. Then
\begin{equation}\label{eq:kbalgebra:point-eval}
 K_b^{(m)}(\zeta)=\prod_{i=0}^{m-1}
 \bigl((1-b_i)t_i+b_i(t_i+1)\bigr).
\end{equation}
\end{lemma}
\begin{proof}
For $b_i=0$ the $i$-th factor on the right is $t_i$, and for $b_i=1$ it is $t_i+1$. Both equal $t_i+b_i=\zeta^{2^i}+b_i$, the corresponding factor in the definition of $K_b^{(m)}(\zeta)$. Multiplication over $i$ gives the equality.
\end{proof}

\begin{definition}[Evaluation curve and scaling]\label{def:kbalgebra:curve}
Fix $\zeta\in\F$. For every $i\in J_m$ put
\[
 t_i(\zeta):=\zeta^{2^i},\qquad
 \gamma_i(\zeta):=2t_i(\zeta)+1,
 \qquad \Gamma_m(\zeta):=\prod_{i\in J_m}\gamma_i(\zeta).
\]
When $\Gamma_m(\zeta)\neq0$, define
\[
 \psi_i(\zeta):=\frac{t_i(\zeta)+1}{\gamma_i(\zeta)},
 \qquad\psi(\zeta):=(\psi_0(\zeta),\ldots,\psi_{m-1}(\zeta)).
\]
\end{definition}

\begin{lemma}[One-factor scaled equality identity]\label{lem:kbalgebra:one-factor-eval}
Suppose $\gamma_i(\zeta)\neq0$. Then for each $b_i\in\{0,1\}$,
\[
 t_i(\zeta)+b_i
 =\gamma_i(\zeta)\bigl((1-b_i)(1-\psi_i(\zeta))+b_i\psi_i(\zeta)\bigr).
\]
\end{lemma}
\begin{proof}
For $b_i=1$, the right side is $\gamma_i\psi_i=t_i+1$ by definition of $\psi_i$. For $b_i=0$, the right side is $\gamma_i(1-\psi_i)=\gamma_i-(t_i+1)=(2t_i+1)-(t_i+1)=t_i$. Both equal $t_i+b_i$. The two Boolean cases exhaust all possibilities.
\end{proof}

\begin{theorem}[Scaled multilinear evaluation]\label{thm:kbalgebra:curve}
For any field $\F$, any table $f:I_m\to\F$, and every $\zeta\in\F$ with $\Gamma_m(\zeta)\neq0$,
\begin{equation}\label{eq:kbalgebra:curve}
 \boxed{U_f^{(m)}(\zeta)=\Gamma_m(\zeta)\,\widetilde f\bigl(\psi(\zeta)\bigr).}
\end{equation}
\end{theorem}
\begin{proof}
Since a product in a field is nonzero precisely when all factors are nonzero, the assumption $\Gamma_m(\zeta)\neq0$ implies $\gamma_i(\zeta)\neq0$ for every $i$. Apply Lemma~\ref{lem:kbalgebra:one-factor-eval} to every factor in $K_b^{(m)}(\zeta)=\prod_i(t_i+b_i)$ to obtain
\[
 K_b^{(m)}(\zeta)=\left(\prod_i\gamma_i(\zeta)\right)
 \prod_i\bigl((1-b_i)(1-\psi_i(\zeta))+b_i\psi_i(\zeta)\bigr)
 =\Gamma_m(\zeta)\lambda_b(\psi(\zeta)).
\]
Multiply this identity by $f(b)$ and sum over $b\in I_m$. The left side is $U_f^{(m)}(\zeta)$ by definition and the right side is $\Gamma_m(\zeta)\widetilde f(\psi(\zeta))$ by Definition~\ref{def:kbalgebra:eq}.
\end{proof}

\begin{proposition}[Exceptional-point bound]\label{prop:kbalgebra:exceptional}
Assume $m\geq1$ and $\operatorname{char}\F\neq2$. Then $\Gamma_m\in\F[X]$ is nonzero of degree $N-1$. Consequently
\[
 \bigl|\{\zeta\in\F:\Gamma_m(\zeta)=0\}\bigr|\leq N-1.
\]
When $\F$ is finite and $|\F|>N-1$, at least $|\F|-(N-1)$ points admit the evaluation formula~\eqref{eq:kbalgebra:curve}.
\end{proposition}
\begin{proof}
Each polynomial $2X^{2^i}+1$ is nonzero with leading coefficient $2\neq0$ and degree $2^i$. Their product $\Gamma_m$ has nonzero leading coefficient $2^m$ and degree $\sum_i2^i=N-1$. A nonzero degree-$d$ polynomial has at most $d$ roots by Lemma~\ref{lem:history:rootcount}. For a finite field, subtracting the number of exceptional roots from the field cardinality gives the final lower bound.
\end{proof}

\begin{corollary}[Characteristic-two evaluation curve]\label{cor:kbalgebra:char2-curve}
If $\operatorname{char}\F=2$, then $\Gamma_m(\zeta)=1$ for every $\zeta$, and
\begin{equation}\label{eq:kbalgebra:char2-curve}
 \boxed{U_f^{(m)}(\zeta)=\widetilde f\bigl(\zeta+1,\zeta^2+1,\ldots,\zeta^{2^{m-1}}+1\bigr).}
\end{equation}
No evaluation points are exceptional.
\end{corollary}
\begin{proof}
In characteristic two, $2t_i+1=1$ for every $i$, so $\Gamma_m=1$ and $\psi_i=(t_i+1)/1=\zeta^{2^i}+1$. Substitute these identities into Theorem~\ref{thm:kbalgebra:curve}. The proof also covers $m=0$ with the empty-product convention.
\end{proof}

\begin{remark}[No division on exceptional points]
When $\Gamma_m(\zeta)=0$ in odd characteristic, the rational expression defining $\psi$ is not available. The universal polynomial identity~\eqref{eq:kbalgebra:point-eval} remains valid at that point; one must use it instead of dividing by a vanishing $\gamma_i$. This distinction is essential for a formalized evaluation algorithm.
\end{remark}

\subsection{Folds on monomials and the general multilinearization functional}\label{sec:kbalgebra:dual}

\begin{definition}[Even--odd polynomial operators]\label{def:kbalgebra:evenodd}
For every polynomial $P(X)=\sum_{j\geq0}p_jX^j\in\F[X]$, define
\[
 E(P)(Y):=\sum_{j\geq0}p_{2j}Y^j,
 \qquad O(P)(Y):=\sum_{j\geq0}p_{2j+1}Y^j.
\]
Thus $P(X)=E(P)(X^2)+XO(P)(X^2)$. Both operators are $\F$-linear. For $r\in\F$ put
\[
 \pfold_r(P):=(2r-1)E(P)+(1-r)O(P).
\]
\end{definition}

\begin{lemma}[Fold of a single monomial]\label{lem:kbalgebra:monomial-fold}
For every integer $j\geq0$ and $r\in\F$,
\[
 \pfold_r(X^{2j})=(2r-1)Y^j,
 \qquad \pfold_r(X^{2j+1})=(1-r)Y^j.
\]
\end{lemma}
\begin{proof}
For $P(X)=X^{2j}$, its even component is $E(P)=Y^j$ and its odd component is $0$ by the coefficient definition; substitute into Definition~\ref{def:kbalgebra:evenodd}. For $P(X)=X^{2j+1}$, the even component is $0$ and the odd component is $Y^j$. Substitute once more.
\end{proof}

\begin{definition}[Complete normalized fold]\label{def:kbalgebra:complete-fold}
For $P\in\F[X]_{<N}$ and $r=(r_0,\ldots,r_{m-1})\in\F^m$, define polynomials $P_0:=P$ and $P_{i+1}:=\pfold_{r_i}(P_i)$ for $0\leq i<m$. The polynomial $P_m$ is constant. Define $\KFold_r(P)\in\F$ to be its constant coefficient. For $m=0$ this is the constant value of $P$.
\end{definition}

\begin{lemma}[Degree halving]\label{lem:kbalgebra:halving}
If $\deg P<2d$, then $\deg\pfold_r(P)<d$, including the zero polynomial. In particular $\KFold_r(P)$ is well defined.
\end{lemma}
\begin{proof}
When all nonzero coefficients of $P$ have indices strictly less than $2d$, both $E(P)$ and $O(P)$ have all coefficients at degrees strictly less than $d$ by Definition~\ref{def:kbalgebra:evenodd}. Their linear combination $\pfold_r(P)$ has no term of degree $\geq d$. Applying this with $d=2^{m-i-1}$ at every step shows that $P_m$ has degree $<1$, i.e. is constant or zero.
\end{proof}

\begin{lemma}[Iterated fold of one monomial]\label{lem:kbalgebra:iterated-monomial}
For $S\subseteq J_m$ and $r\in\F^m$,
\begin{equation}\label{eq:kbalgebra:monomial-fold}
 \boxed{\KFold_r(X^{e(S)})=
 \left(\prod_{i\in J_m\setminus S}(2r_i-1)\right)
 \left(\prod_{i\in S}(1-r_i)\right).}
\end{equation}
\end{lemma}
\begin{proof}
We induct on $m$. At $m=0$, $S=\varnothing$, $e(S)=0$ and the fold is $X^0=1$, while both products are empty and equal $1$. Suppose $m\geq1$. Write $e(S)=b_0+2j$, where $b_0=1$ if $0\in S$ and $0$ otherwise, and $j=\sum_{i\in S,i>0}2^{i-1}$. By Lemma~\ref{lem:kbalgebra:monomial-fold}, the first fold yields $(2r_0-1)Y^j$ if $b_0=0$, or $(1-r_0)Y^j$ if $b_0=1$. All remaining folds are $\F$-linear, so this scalar factors out. Apply the induction hypothesis to $Y^j$ and the remaining challenge coordinates $(r_1,\ldots,r_{m-1})$. Multiplying the scalar from the first fold by the inductively obtained factors gives exactly the two products in~\eqref{eq:kbalgebra:monomial-fold}.
\end{proof}

\begin{definition}[Binary-exponent multilinearization]\label{def:kbalgebra:multilinearization}
Let $\mathscr V_m:=\F[Z_0,\ldots,Z_{m-1}]_{\mathrm{ml}}$ be the $\F$-linear span of squarefree monomials. Define the unique $\F$-linear map $\mathscr M_m:\F[X]_{<N}\to\mathscr V_m$ by
\[
 \mathscr M_m(X^{e(S)})=\prod_{i\in S}Z_i.
\]
For $P(X)=\sum_{S\subseteq J_m}a_SX^{e(S)}$ we consequently have $\mathscr M_m(P)(Z)=\sum_S a_S\prod_{i\in S}Z_i$. No multiplicativity of $\mathscr M_m$ is assumed.
\end{definition}

\begin{lemma}[Multilinearization is a vector-space isomorphism]\label{lem:kbalgebra:multilinear-iso}
The map $\mathscr M_m$ is an $\F$-linear isomorphism $\F[X]_{<N}\to\mathscr V_m$.
\end{lemma}
\begin{proof}
By Lemma~\ref{lem:kbalgebra:binary-bijection}, the monomials $X^{e(S)}$ for $S\subseteq J_m$ form the usual basis $1,X,\ldots,X^{N-1}$. The squarefree monomials $\prod_{i\in S}Z_i$ form an $N$-element basis of $\mathscr V_m$ by its definition. The map $\mathscr M_m$ sends the first basis bijectively onto the second; a linear map doing so has a linear inverse determined by the reverse basis correspondence.
\end{proof}

\begin{lemma}[Multilinearization of kernels, every characteristic]\label{lem:kbalgebra:multilinear-kernel}
For every $b\in I_m$ and every field $\F$,
\begin{equation}\label{eq:kbalgebra:ml-kernel}
 \boxed{\mathscr M_m(K_b^{(m)})(Z)=\prod_{i\in J_m}(Z_i+b_i).}
\end{equation}
\end{lemma}
\begin{proof}
By Lemma~\ref{lem:kbalgebra:product-expansion},
$K_b^{(m)}=\sum_{S\subseteq J_m}(\prod_{i\notin S}b_i)X^{e(S)}$.
Apply the $\F$-linearity of $\mathscr M_m$ and then its defining formula on each monomial:
\[
 \mathscr M_m(K_b^{(m)})(Z)=\sum_{S\subseteq J_m}
 \left(\prod_{i\notin S}b_i\right)\left(\prod_{i\in S}Z_i\right).
\]
By distributivity, the right-hand side is precisely the expansion of $\prod_i(Z_i+b_i)$. This proof explicitly expands the product; it does not use, or assert, that $\mathscr M_m$ is a ring homomorphism.
\end{proof}

\begin{definition}[Tensor evaluation functional]\label{def:kbalgebra:weighted-eval}
For $r\in\F^m$ define an $\F$-linear map $\Lambda_r:\mathscr V_m\to\F$ on squarefree monomials by
\[
 \Lambda_r\left(\prod_{i\in S}Z_i\right):=
 \prod_{i\notin S}(2r_i-1)\prod_{i\in S}(1-r_i).
\]
This is well defined because the squarefree monomials form a basis.
\end{definition}

\begin{theorem}[Complete fold as a tensor functional]\label{thm:kbalgebra:tensor-fold}
For every field $\F$, every $m\geq0$, $P\in\F[X]_{<N}$, and $r\in\F^m$,
\begin{equation}\label{eq:kbalgebra:tensor-fold}
 \boxed{\KFold_r(P)=\Lambda_r(\mathscr M_m(P)).}
\end{equation}
\end{theorem}
\begin{proof}
Both sides define $\F$-linear functionals in $P$: the left side is a composition of $\F$-linear folds by Definition~\ref{def:kbalgebra:evenodd}, and the right side is a composition of the $\F$-linear maps $\mathscr M_m$ and $\Lambda_r$. It suffices to check equality on the basis monomials $X^{e(S)}$. The left side on such a monomial is given explicitly in Lemma~\ref{lem:kbalgebra:iterated-monomial}. The right side is $\Lambda_r(\prod_{i\in S}Z_i)$ by Definition~\ref{def:kbalgebra:multilinearization}, which equals the same product by Definition~\ref{def:kbalgebra:weighted-eval}. Thus the functionals agree on a basis and therefore on every $P$.
\end{proof}

\begin{corollary}[Ordinary evaluation when all weights are invertible]\label{cor:kbalgebra:weighted-eval}
If $2r_i-1\neq0$ for every $i$, then
\begin{equation}\label{eq:kbalgebra:weighted-eval}
 \KFold_r(P)=\left(\prod_i(2r_i-1)\right)
 \mathscr M_m(P)\left(\frac{1-r_0}{2r_0-1},\ldots,\frac{1-r_{m-1}}{2r_{m-1}-1}\right).
\end{equation}
\end{corollary}
\begin{proof}
Apply Theorem~\ref{thm:kbalgebra:tensor-fold} and expand $P$ in monomials. In the summand indexed by $S$, each factor with $i\in S$ is $(1-r_i)$; since $2r_i-1$ is invertible, this equals $(2r_i-1)\cdot(1-r_i)/(2r_i-1)$. Factoring $\prod_i(2r_i-1)$ from all summands leaves the evaluation of $\mathscr M_m(P)$ at the displayed point.
\end{proof}

\begin{corollary}[Characteristic-two fold identity]\label{cor:kbalgebra:char2-multilinear}
If $\operatorname{char}\F=2$, then $2r_i-1=1$ and $1-r_i=1+r_i$ for every $i$, and hence
\begin{equation}\label{eq:kbalgebra:char2-multilinear}
 \boxed{\KFold_r(P)=\mathscr M_m(P)(1+r_0,\ldots,1+r_{m-1}).}
\end{equation}
\end{corollary}
\begin{proof}
In characteristic two the hypotheses of Corollary~\ref{cor:kbalgebra:weighted-eval} hold automatically because $2r_i-1=-1=1$. Substituting $2r_i-1=1$ and $1-r_i=1+r_i$ into~\eqref{eq:kbalgebra:weighted-eval} yields the statement.
\end{proof}

\begin{remark}[Critical characteristic distinction]
The identity~\eqref{eq:kbalgebra:char2-multilinear} is special to characteristic two. In odd characteristic the correct universal identity is the weighted tensor functional~\eqref{eq:kbalgebra:tensor-fold}; one must not evaluate $\mathscr M_m(P)$ at $1+r$ without the appropriate weights and reparametrization.
\end{remark}

\subsection{Classical FRI, field extension, and degree-aware masking}\label{sec:kbalgebra:fri}

\begin{definition}[Classical even--odd FRI fold]\label{def:kbalgebra:classical}
For $P\in\F[X]$ and $\theta\in\F$, put
\[
 \mathsf{cfold}_\theta(P):=E(P)+\theta O(P),
\]
where $E$ and $O$ are the operators of Definition~\ref{def:kbalgebra:evenodd}.
\end{definition}

\begin{lemma}[Classical fold of the kernel coordinate pair]\label{lem:kbalgebra:classical-coordinates}
Let $m\geq1$ and write a table $f:I_m\to\F$ as $f_0(b'):=f(0,b')$ and $f_1(b'):=f(1,b')$. Then
\[
 \mathsf{cfold}_\theta(U_f^{(m)})
 =U_{\theta f_0+(1+\theta)f_1}^{(m-1)}.
\]
\end{lemma}
\begin{proof}
By Corollary~\ref{cor:kbalgebra:recursion},
\[
 U_f^{(m)}(X)=XU_{f_0}^{(m-1)}(X^2)+(X+1)U_{f_1}^{(m-1)}(X^2).
\]
Collecting the even and odd terms gives $E(U_f)=U_{f_1}^{(m-1)}$ and $O(U_f)=U_{f_0}^{(m-1)}+U_{f_1}^{(m-1)}$. Thus
\[
 \mathsf{cfold}_\theta(U_f)=U_{f_1}+\theta(U_{f_0}+U_{f_1})
 =U_{\theta f_0+(1+\theta)f_1}
\]
by linearity of $f\mapsto U_f$.
\end{proof}

\begin{lemma}[Classical versus normalized fold]\label{lem:kbalgebra:fold-reparam}
Let $\theta\in\F$ satisfy $1+2\theta\neq0$, and set
\[
 r:=\frac{1+\theta}{1+2\theta}.
\]
Then for every polynomial $P\in\F[X]$,
\begin{equation}\label{eq:kbalgebra:fold-reparam}
 \boxed{\mathsf{cfold}_\theta(P)=(1+2\theta)\pfold_r(P).}
\end{equation}
Conversely, for $r\in\F$ with $2r-1\neq0$, setting $\theta:=(1-r)/(2r-1)$ gives
\[
 \pfold_r(P)=(2r-1)\mathsf{cfold}_\theta(P).
\]
\end{lemma}
\begin{proof}
For the first identity,
\[
 2r-1=\frac{2+2\theta-(1+2\theta)}{1+2\theta}
 =\frac1{1+2\theta},
 \qquad 1-r=\frac{\theta}{1+2\theta}.
\]
Therefore $(1+2\theta)((2r-1)E(P)+(1-r)O(P))=E(P)+\theta O(P)=\mathsf{cfold}_\theta(P)$. For the converse, substituting $\theta=(1-r)/(2r-1)$ into $(2r-1)(E+\theta O)$ gives $(2r-1)E+(1-r)O=\pfold_r(P)$. Every division is legal by the explicitly stated nonvanishing hypotheses.
\end{proof}

\begin{lemma}[One classical fold preserves evaluation on a square]\label{lem:kbalgebra:classical-eval}
For every $P\in\F[X]$, every $\zeta\in\F$, and $\theta=\zeta$,
\[
 P(\zeta)=\mathsf{cfold}_{\zeta}(P)(\zeta^2).
\]
\end{lemma}
\begin{proof}
By Definition~\ref{def:kbalgebra:evenodd}, $P(\zeta)=E(P)(\zeta^2)+\zeta O(P)(\zeta^2)$. The right side is exactly $\mathsf{cfold}_\zeta(P)(\zeta^2)$ by Definition~\ref{def:kbalgebra:classical}.
\end{proof}

\begin{theorem}[Classical folding computes univariate evaluation]\label{thm:kbalgebra:classical-evaluation}
Fix $\zeta\in\F$ and $P\in\F[X]_{<N}$. Define $P_0=P$ and $P_{i+1}:=\mathsf{cfold}_{\zeta^{2^i}}(P_i)$ for $0\leq i<m$. Then $P_m$ is constant and
\[
 \boxed{P_m=P(\zeta).}
\]
\end{theorem}
\begin{proof}
The classical fold maps polynomials of degree less than $2d$ to polynomials of degree less than $d$, because $E$ and $O$ have that degree bound by Definition~\ref{def:kbalgebra:evenodd}. Therefore $P_m$ is constant. Applying Lemma~\ref{lem:kbalgebra:classical-eval} with $P_i$ and evaluation point $\zeta^{2^i}$ shows
\[
 P_i(\zeta^{2^i})=P_{i+1}(\zeta^{2^{i+1}})
\]
for every $0\leq i<m$. Chaining these $m$ equalities yields $P(\zeta)=P_m(\zeta^{2^m})=P_m$, since $P_m$ is constant.
\end{proof}

\begin{lemma}[Naturality under field extensions]\label{lem:kbalgebra:naturality}
Let $\iota:\F\hookrightarrow\mathbb K$ be an injective field homomorphism. Extend it coefficientwise to polynomials and pointwise to tables. Then
\[
 \iota(K_b^{(m)})=K_b^{(m)},\qquad
 \iota(U_f^{(m)})=U_{\iota\circ f}^{(m)},
\]
and for every $P\in\F[X]$ and $r\in\F$,
\[
 \iota(\pfold_r(P))=\pfold_{\iota(r)}(\iota(P)).
\]
Consequently the kernel-coordinate transforms and complete normalized folds commute with $\iota$.
\end{lemma}
\begin{proof}
Every coefficient of $K_b^{(m)}$ is an integer expression in $0$ and $1$, fixed by the homomorphism $\iota$, proving the first equality. Apply $\iota$ to the finite sum defining $U_f$ and use preservation of addition and multiplication to obtain the second. Even--odd coefficient extraction commutes with $\iota$ because $\iota$ acts coefficientwise: $E(\iota(P))=\iota(E(P))$ and $O(\iota(P))=\iota(O(P))$. Applying $\iota$ to $(2r-1)E(P)+(1-r)O(P)$ therefore gives the third equality. Repeating it round by round proves the corresponding complete fold identity. Uniqueness of kernel coordinates from Theorem~\ref{thm:kbalgebra:basis} gives the transform compatibility.
\end{proof}

\begin{lemma}[Degree of a high-degree masked kernel polynomial]\label{lem:kbalgebra:mask-degree}
Let $m\geq1$, $N=2^m$, $f:I_m\to\F$, and $R\in\F[X]$ be nonzero of degree $d\geq0$. Then
\begin{equation}\label{eq:kbalgebra:mask-degree}
 \boxed{\deg(U_f^{(m)}+X^N R)=N+d.}
\end{equation}
\end{lemma}
\begin{proof}
Corollary~\ref{cor:kbalgebra:monic} gives $\deg U_f^{(m)}\leq N-1$, unless the polynomial vanishes. Because $R$ is nonzero of degree $d$, the polynomial $X^NR$ has degree exactly $N+d$ with the same nonzero leading coefficient as $R$. There is no term of degree $N+d$ in $U_f$, so the top coefficient of $X^NR$ cannot cancel in the sum. The sum has degree $N+d$.
\end{proof}

\begin{corollary}[The code dimension must change]\label{cor:kbalgebra:mask-rate}
Suppose $\mathcal D\subseteq\F$ has $M$ distinct elements and $M\geq N+d+1$, with $R\neq0$ of degree $d$. Then the evaluation word of $U_f+X^NR$ does not lie in $\operatorname{RS}_{\F}(\mathcal D,N)$. It belongs to $\operatorname{RS}_{\F}(\mathcal D,N+d+1)$.
\end{corollary}
\begin{proof}
The second claim follows from Lemma~\ref{lem:kbalgebra:mask-degree}: the polynomial has degree $N+d<N+d+1$. Suppose its evaluation word also lay in $\operatorname{RS}_{\F}(\mathcal D,N)$. Then some polynomial $T$ of degree less than $N$ would agree with $U_f+X^NR$ at every point of $\mathcal D$. Their difference is nonzero of degree $N+d$, because the leading term of $X^NR$ cannot cancel with terms of degree less than $N$. This difference has at least $M\geq N+d+1$ distinct roots, contradicting Lemma~\ref{lem:history:rootcount}. Hence the word is outside the original degree-$N$ code.
\end{proof}

\subsection{Kernel-mask subspaces and bounded-query hiding}\label{sec:kbalgebra:maskquery}
Throughout this subsection, $\F$ is a \emph{finite} field.

\begin{definition}[Random low-degree kernel mask]\label{def:kbalgebra:random-mask}
Let $\F$ be a finite field. Fix $0\leq t\leq m$ and let $B=2^t$. Let $\mathcal H_t\subseteq\F^{I_m}$ be the coefficient subspace of Corollary~\ref{cor:kbalgebra:dimension}. Sample a random table $H$ uniformly from $\mathcal H_t$, and let $R:=U_H^{(m)}$ be its encoded polynomial, which satisfies $\deg R<B$. For an arbitrary fixed $P\in\F[X]$, the 
\emph{high-degree masked polynomial} is
\[
 Q(X):=P(X)+X^N R(X),\qquad N=2^m.
\]
\end{definition}

\begin{lemma}[Uniformity of the mask polynomial]\label{lem:kbalgebra:uniform-polynomial}
Under Definition~\ref{def:kbalgebra:random-mask}, the random polynomial $R$ is uniform on the $B$-dimensional vector space $\F[X]_{<B}$.
\end{lemma}
\begin{proof}
Theorem~\ref{thm:kbalgebra:low-degree} identifies $\mathcal H_t$ with the tables $\varphi:I_t\to\F$ via an injective linear map onto $\mathcal H_t$. The same theorem identifies $U_H^{(m)}$ with $U_\varphi^{(t)}$. By the basis theorem (Theorem~\ref{thm:kbalgebra:basis}) with $t$ variables, the map $\varphi\mapsto U_\varphi^{(t)}$ is a linear bijection from $\F^{I_t}$ to $\F[X]_{<B}$. The composition $H\mapsto R$ is therefore a bijection between the finite sets $\mathcal H_t$ and $\F[X]_{<B}$. A bijection sends the uniform distribution to the uniform distribution.
\end{proof}

\begin{lemma}[Uniform answers at prescribed nonzero points]\label{lem:kbalgebra:uniform-queries}
Let $1\leq q\leq B$, and let $\xi_1,\ldots,\xi_q\in\F^\times$ be pairwise distinct. With $R$ as in Definition~\ref{def:kbalgebra:random-mask}, for any fixed $P\in\F[X]$ the random vector
\[
 (Q(\xi_1),\ldots,Q(\xi_q))\in\F^q,
 \qquad Q=P+X^N R,
\]
is uniform on $\F^q$ and independent of the choice of $P$.
\end{lemma}
\begin{proof}
By Lemma~\ref{lem:kbalgebra:uniform-polynomial}, $R$ is uniform on $\F[X]_{<B}$. Let $\operatorname{Ev}_q:\F[X]_{<B}\to\F^q$ be evaluation at the $q$ specified points. Since $q\leq B$, the Lagrange polynomials of Lemma~\ref{lem:security:lagrange} have degree $<q\leq B$, and prove that $\operatorname{Ev}_q$ is surjective. As an $\F$-linear surjection between finite vector spaces, all its fibers are cosets of the same kernel and have equal cardinality. Consequently the evaluation vector $(R(\xi_1),\ldots,R(\xi_q))$ is uniform on $\F^q$. Every $\xi_j$ is nonzero, so multiplication of coordinate $j$ by $\xi_j^N$ is invertible. The diagonal map $(a_j)_j\mapsto(\xi_j^Na_j)_j$ is a bijection of $\F^q$, and adding the fixed vector $(P(\xi_j))_j$ is another bijection. Their composition sends the uniform vector of $R$-evaluations to a uniform vector of $Q$-evaluations, independently of $P$.
\end{proof}

\begin{theorem}[Adaptive bounded-query hiding for one masked polynomial]\label{thm:kbalgebra:adaptive-queries}
Let $R$ be sampled as in Definition~\ref{def:kbalgebra:random-mask} with $B=2^t$. Consider a possibly randomized verifier whose random tape is independent of $R$ and which adaptively queries the oracle $Q=P+X^N R$ at at most $B$ distinct nonzero field points, with repeat queries allowed but returning their previous value. The strategy for selecting the next query depends only on public parameters, the verifier's tape, and the earlier query-and-answer pairs. If the verifier receives no other $P$-dependent messages, the distribution of its complete query-and-answer record is independent of the fixed secret polynomial $P$.
\end{theorem}
\begin{proof}
Condition first on a fixed random tape of the verifier, so the query strategy is deterministic. We prove by induction on the number $q$ of distinct queried points that the conditional law of the next \emph{fresh} answer, given any realized history of $q<B$ distinct queries, is uniform on $\F$ and independent of $P$. The base $q=0$ is Lemma~\ref{lem:kbalgebra:uniform-queries} with one point. For the induction step, fix a history of $q$ pairwise distinct nonzero indices $\xi_1,\ldots,\xi_q$ and answers $a_1,\ldots,a_q$ with positive probability. The verifier's deterministic strategy fixes its next fresh point $\xi_{q+1}$ from this history. Since it is distinct from previous points and $q+1\leq B$, Lemma~\ref{lem:kbalgebra:uniform-queries} applied to this fixed $(q+1)$-tuple states that the joint vector of answers is uniform on $\F^{q+1}$. Thus
\[
 \Pr[Q(\xi_{q+1})=a\mid Q(\xi_j)=a_j\ (1\leq j\leq q)]
 =\frac{|\F|^{-(q+1)}}{|\F|^{-q}}=|\F|^{-1},
\]
for each $a\in\F$, independently of $P$. On a repeated query the answer is deterministically the previously observed value, again independent of $P$ conditional on the history. By induction on query steps, the distributions of the entire records coincide for all choices of $P$. Finally average over the verifier's independent random tape; the mixture is the same for every $P$.
\end{proof}

\begin{theorem}[Mask-query and code-rate tradeoff]\label{thm:kbalgebra:mask-rate-tradeoff}
Let $\F$ be a finite field, let $m\geq1$, let $N=2^m$, and let $0\leq t\leq m$ with $B=2^t$. Let $\mathcal D\subseteq\F^\times$ have $M\geq N+B$ pairwise distinct elements. For a fixed secret table $f:I_m\to\F$, sample $R$ uniformly from $\F[X]_{<B}$ and set
\[
 Q(X):=U_f^{(m)}(X)+X^N R(X).
\]
Then the following three statements hold:
\begin{enumerate}[label=\textup{(\roman*)}]
\item $\operatorname{ev}_{\mathcal D}(Q)\in\operatorname{RS}_{\F}(\mathcal D,N+B)$ for every $R$.
\item A verifier making at most $B$ adaptive distinct queries to this one oracle, at nonzero domain points and without receiving other witness-dependent messages, obtains a query-and-answer distribution independent of $f$.
\item The code $\operatorname{RS}_{\F}(\mathcal D,N+B)$ has minimum distance at least $M-(N+B)+1$, and a received word can be within Hamming distance $\tau$ of at most one codeword whenever $2\tau<M-(N+B)+1$.
\end{enumerate}
\end{theorem}
\begin{proof}
For (i), $\deg U_f^{(m)}\leq N-1$ and $\deg R<B$. Therefore $\deg(X^N R)<N+B$, with the zero polynomial also allowed. The same degree bound holds for their sum $Q$, which by definition belongs to the evaluation code of dimension bound $N+B$. For (ii), Lemma~\ref{lem:kbalgebra:uniform-polynomial} shows that uniform kernel masks from $\mathcal H_t$ are equivalent to uniform polynomials of degree $<B$, so Theorem~\ref{thm:kbalgebra:adaptive-queries} applies verbatim to the oracle $Q$ on the nonzero domain $\mathcal D$. For (iii), two distinct polynomials of degree $<N+B$ differ by a nonzero polynomial of degree at most $N+B-1$. The root-count lemma (Lemma~\ref{lem:history:rootcount}) shows that their evaluation words agree in at most $N+B-1$ of the $M$ positions. Thus their Hamming distance is at least $M-(N+B)+1$. If two codewords both had distance at most $\tau$ from one received word, the Hamming triangle inequality (Lemma~\ref{lem:security:hamming-triangle}) would bound their mutual distance by $2\tau$, contrary to the strict hypothesis.
\end{proof}

\begin{remark}[Exact scope of the bounded-query theorem]\label{rem:kbalgebra:mask-scope}
Theorem~\ref{thm:kbalgebra:adaptive-queries} concerns \emph{only} direct oracle evaluations of a single high-degree masked polynomial at nonzero points. It does not prove zero knowledge of KBFold: the original protocol reveals additional sumcheck round polynomials, correlated folded words, a last-round table, and authenticated commitment openings. Those additional messages can invalidate the theorem's hypotheses. Moreover, to use the masked word in a proximity protocol the Reed--Solomon degree and final closure must be modified as proved in Corollary~\ref{cor:kbalgebra:mask-rate} and Theorem~\ref{thm:security:naive-closure}.
\end{remark}


\subsection{Formal status and next algebraic dependencies}\label{sec:kbalgebra:status}
The foundations above depend only on field arithmetic, polynomial coefficients, finite sums, and finite-dimensional linear algebra, except for the explicitly finite-field masking statements in Section~\ref{sec:kbalgebra:maskquery}. The tensor fold in Theorem~\ref{thm:kbalgebra:tensor-fold} is valid in every characteristic; the ordinary evaluation formula of Corollary~\ref{cor:kbalgebra:char2-multilinear} is the characteristic-two specialization. These are not interchangeable in the odd-characteristic KBFold protocol. In the next chapter, the typed mathematical statements on multilinearization, folding, and masks can be imported as standalone lemmas rather than re-proved inside a cryptographic security argument. The dedicated theorem ledger accompanying the source provides a proposed Lean name for each declaration. No new Lean machine-checked proof is asserted here.
